\subsection{Restricting to a linear subcode}
\label{sec:retained-stability}
The atomic inequality remains valid when only some transversals are
retained. The support-stability statement needs a little more care,
since its selected full-code transversal might have been omitted.
Let $\mathcal A$ be any nonempty collection of retained transversals,
let $Q_{\mathcal A}$ denote the corresponding weighted coefficient sum,
and let $\mathcal D_{\mathcal A,\mathrm{odd}}$ be their odd downward
Fourier space. For $\ell$ rail pairs, one has the safe estimate
\begin{equation}\label{eq:retained-stability}
 \|\phi_u-P_{\mathcal D_{\mathcal A,\mathrm{odd}}}\phi_u\|_2^2
 \le 2^{20}\ell^2R^2\|u\|_2
       \bigl(\|u\|_2-Q_{\mathcal A}(\phi_u)\bigr).
\end{equation}

Here is a proof which does not assume that a balanced hardtanh ray
has the optimal gain. Fix one rail $B$. In the positive absolute
integral, sum all full-code terms which use $B$. For balanced rails,
apply the shared-factor Cauchy inequality and the unpaired-block
argument of Lemma~\ref{lem:random-min}, treating $B$ itself as a fixed
selected unpaired block. Discard the cosine factor of its opposite
rail. Taking the finite derivative in a coordinate of $B$ bounds each
of the two unweighted shared-factor integrals by
$\kappa_{D'}^{\rm R}a_B$, where $a_B$ is its common amplitude and
$D'\ge\ell R+2$. This argument does not require that $B$ have the
larger amplitude in its pair. Since
$w_{\max}\le2(2\ell R+1)$ and
$\kappa_{D'}^{\rm R}\le1/\sqrt{D'}$, its weighted hardtanh
marginal is at most $8\sqrt{\ell R}\,a_B$ (the finitely small
endpoint cases are also covered by $R\ge12$).
Block H\"older applied to this marginal, followed by the geometric
mean--quadratic mean inequality in the sampled coordinate of $B$,
therefore bounds the arbitrary-coordinate marginal by
\[
 8\sqrt{\ell R}\left(\frac1{|B|}\sum_{j\in B}u_j^2\right)^{1/2}
 \le8\sqrt\ell\,\|u_B\|_2.
\]
Consequently the positive potential over terms using at least one
rail from a specified collection of one rail per pair is at most
$8\ell$ times the Euclidean norm of all those rail weights.

Normalize $\|u\|_2=1$ and write
$\epsilon=1-Q_{\mathcal A}(\phi_u)$. The full-code support argument
in the preceding subsection applies because the retained coefficient
sum is bounded by the full positive potential. It selects one full
transversal, and, after removing bias and irrelevant coordinates,
its discarded parameter norm $\rho$ satisfies
$\rho^2\le2^8R^2\epsilon$. If that transversal is retained, the
one-Lipschitz argument already proves \eqref{eq:retained-stability}.
Otherwise every retained transversal uses at least one of these
unused rails. The marginal estimate gives
\[
 1-\epsilon=Q_{\mathcal A}(\phi_u)
 \le8\ell\rho\le2^7\ell R\sqrt\epsilon.
\]
For $\epsilon<1/2$ this implies
$\epsilon\ge1/(2^{16}\ell^2R^2)$; for $\epsilon\ge1/2$ the same
conclusion is immediate. Approximate by the zero function, which
belongs to every retained downward space, and use
$\E\phi_u^2\le\E(u\cdot\bar x)^2=1$. This proves the normalized
bound. Rescaling the parameter direction, with the auxiliary
hardtanh threshold rescaled inversely, gives the stated homogeneous
parameter bound; equivalently repeat the same argument after dividing
all positive potentials and parameter norms by $\|u\|_2$.
